\section{Kritische Würdigung und Gestaltungsempfehlungen}\label{sec:wuerdigung}
Das Muster aus Kapitel~\ref{subsec:produktivitaet} setzt sich fort: Die Befunde unterscheiden sich darin, worin die Werkzeuge eingebettet sind. Der DORA-Report zieht aus seinen eigenen gegenläufigen Zahlen den Schluss, dass ein verbesserter Entwicklungsprozess die Auslieferung nicht von selbst verbessert, solange Grundlagen wie kleine Änderungspakete und belastbare Testmechanismen fehlen.\footcite[\vglf][\pagef 40]{DORA.2024} Die Mehrfallstudie stützt das von der Teamseite: Die Werkzeuge wirken als persönliche Assistenz, ohne die Zusammenarbeit zu verändern.\footcite[\vglf][\pagef 299]{Neumann.2026} Daraus ergibt sich die These dieser Arbeit: Über Nutzen und Schaden entscheidet nicht das Werkzeug, sondern die Praktik, die es umgibt. Review, Definition of Done und Retrospektive sind zugleich Verstärker der Potenziale und Schutz vor den Risiken.

Fünf Empfehlungen folgen daraus. Erstens ist die Definition of Done um Kriterien für generierten Code zu erweitern; angesichts der Verwundbarkeitsquoten aus Kapitel~\ref{subsec:qualitaetsrisiken} gehört eine automatisierte Sicherheitsprüfung zur Mindestanforderung, und im Review ist generierter Code als solcher auszuweisen. Zweitens sind Nutzungsrichtlinien verbindlich zu fassen, da Beschäftigte sonst nach eigenem Ermessen entscheiden.\footcite[\vglf][\pagef 298]{Neumann.2026} Drittens ist der Einsatz nach Erfahrungsstand zu differenzieren: Weil die Gewinne bei Berufseinsteigern auftreten und bei erfahrenen Entwicklern ausbleiben oder ins Gegenteil umschlagen,\footcites(\vglf)()[][\pagef 5]{Gambacorta.2024}[][\pagef 2]{Becker.2025} liegt deren Beitrag eher in Architektur und Review als in generierungsnaher Arbeit. Viertens erhalten drei von sechs beobachteten Interaktionsmustern das Lernen auch unter KI-Nutzung;\footcite[\vglf][\pagef 1]{Shen.2026} maßgeblich ist demnach die kognitive Beteiligung, nicht der Verzicht. Fünftens lässt sich der Wahrnehmungslücke aus Kapitel~\ref{subsec:kompetenzrisiken} mit dem vorhandenen Instrumentarium begegnen: Die Retrospektive ist die Praktik, die die Selbsteinschätzung korrigiert -- allerdings nur, wenn ihr gemessene Durchlaufzeiten zugrunde liegen und nicht Eindrücke. Das entspricht dem Zweck der Größe: Team Velocity dient realistischen Vorhersagen, nicht der Maximierung der Arbeitsgeschwindigkeit.\footcite[\vglf][\pagef 133]{Preussig.2020}

Die Grenzen der Arbeit sind zu benennen. Die Literaturbasis ist jung und teils vorläufig: Zwei tragende Studien liegen als Vorabfassung vor, eine stammt vom Anbieter selbst, eine weitere umfasst sechzehn Personen. Laborbefunde und Selbstauskünfte lassen sich nicht ohne Weiteres auf agile Teams übertragen, und der Ursachenzusammenhang hinter den Prozesskennzahlen ist im Report selbst nur eine Vermutung. Die Empfehlungen sind daher begründete Ableitungen, keine gesicherten Wirkungsaussagen.
